\section{Conclusion}\label{sec:conclusion120}
Neural Bellman Operators connect learned continuation operators to feasible
decisions and the economic value of their implemented future behavior. The
construction requires a fixed original recursion and a complete accounting of
the information used by its critic, actor and numerical certificate. The
constrained investment model makes these requirements explicit while retaining
continuous innovations and nonlinear recursive risk.

The new action-comparison theory uses the geometry of the original Bellman
problem. A constrained state-gradient modulus yields an action-scaled
continuation-difference bound; convex integration and a constrained risk-defect
transfer then provide computable comparison intervals. Qualification and a
strict original-policy improvement bound jointly govern replacement of a
completed incumbent. The proof also identifies what additional precision can
resolve and what remains hidden behind a uniform future-policy allowance. This
distinction gives a concrete stopping criterion and a specified information
problem for subsequent NBO algorithms.

R120-INTEGRATION-PENDING: integrate the actual original-policy, population-gain,
all-comparator and complete-cost findings here after the declared execution.

The remaining research priorities concern the joint design of continuation
learning and paid economic decisions. They include acquisition losses that
reduce action-comparison work, information about the actual future controller
that resolves gains below the present allowance, and matched experiments with
longer horizons, multiple controls and higher-dimensional continuous risk.
The retained stochastic-control, recursive-utility and stopped-boundary results
provide separate rigorous foundations for this program. Progress on these
problems must preserve the original objective, give conventional methods the
same usable information, and demonstrate the economic value of the resulting
policies together with their full computational cost.
